% !TeX root = ../凸优化笔记.tex
% 由主文档通过 \input 载入；请编译主文档。
\setcounter{chapter}{0} % 保留原书章号。
\chapter{数学基础与符号补充}
\label{app:co-background}

本附录集中列出正文使用的基础工具。首次阅读时可以按引用查阅；各结论的条件也是分析优化模型时应保留的条件。

\section{范数、内积与对偶范数}

范数 $\|x\|$ 满足正定性、绝对齐次性 $\|\alpha x\|=|\alpha|\|x\|$ 和三角不等式。对 $1\leq p<\infty$，$\|x\|_p=(\sum_i|x_i|^p)^{1/p}$；$\|x\|_\infty=\max_i|x_i|$。若 $p,q>1$ 且 $1/p+1/q=1$，Young 不等式 $ab\leq a^p/p+b^q/q$ 对非负 $a,b$ 成立，它可由固定 $a$ 后求右侧减左侧关于 $b$ 的最小值证明。对归一化分量应用并求和，即得 Hölder 不等式
\[
 |x^{\mathsf T}y|\leq\|x\|_p\|y\|_q.
\]
端点形式为 $|x^{\mathsf T}y|\leq\|x\|_1\|y\|_\infty$；$p=q=2$ 为柯西--施瓦茨不等式。

对偶范数 $\|y\|_*=\max_{\|x\|\leq1}y^{\mathsf T}x$ 的最大值在有限维紧单位球上取到。Hölder 不等式给出上界，选取同号且大小成适当幂次的分量可取等号，所以 $\ell_p$ 的对偶为 $\ell_q$。$P\succ0$ 时，$\sqrt{x^{\mathsf T}Px}$ 的对偶范数为 $\sqrt{y^{\mathsf T}P^{-1}y}$。

\section{有限维拓扑与极值存在性}

有限维空间中，集合紧当且仅当闭且有界；紧集中的任意序列都有收敛于该集合的子列。连续函数在非空紧集上取到最大、最小值。较弱的下半连续性要求 $f(x)\leq\liminf_k f(x_k)$ 对所有 $x_k\to x$ 成立；在非空紧集上，它已足以保证有限值函数取到最小值。

\begin{proposition}
设 $f:C\to\mathbb R$ 下半连续，存在实数 $\alpha$ 使次水平集 $S_\alpha=\{x\in C:f(x)\leq\alpha\}$ 非空且紧，则 $f$ 在 $C$ 上有全局最小点。
\end{proposition}
\begin{proof}
先在紧集 $S_\alpha$ 上取极小化序列。若其函数值无下界，紧性给出收敛子列，下半连续性将迫使有限函数值不大于 $-\infty$，矛盾。因此下确界有限，取收敛子列后，下半连续性说明极限点取到该下确界。$S_\alpha$ 外的函数值大于 $\alpha$，故该点也是全局最小点。
\end{proof}

强制性（coercivity）指 $\|x\|_2\to\infty$ 时 $f(x)\to+\infty$。在非空闭集上的连续强制函数有非空紧次水平集，因而取到最小值。开集上的连续函数则可能在缺失的边界逼近下确界，如 $f(x)=x$ 在 $(0,1)$ 上。

\begin{example}[用次水平集证明最优解存在]
设 $C\subseteq\mathbb R^n$ 非空闭，$A,b$ 给定，$\alpha>0$。证明 $f(x)=\|Ax-b\|_2^2/2+\alpha\|x\|_2^2/2$ 在 $C$ 上达到最小值。
\end{example}
\begin{solve}
取任意 $x_0\in C$。在次水平集 $f(x)\leq f(x_0)$ 中，由 $f(x)\geq\alpha\|x\|_2^2/2$ 得 $\|x\|_2\leq\sqrt{2f(x_0)/\alpha}$。该次水平集非空、闭且有界，故紧，连续目标在其上达到最小值。若 $C$ 还凸，强凸性进一步给出唯一性；没有凸性时，上面的存在性证明仍然成立，但唯一性不再有保证。
\end{solve}

\section{微分与矩阵分析}

可微性表示 $f(x+h)=f(x)+\nabla f(x)^{\mathsf T}h+o(\|h\|_2)$，要求余项在所有趋零方向上一致地小。沿每条直线有方向导数并不足够：令 $f(0,0)=0$，非原点处 $f(x,y)=x^3/(x^2+y^2)$，其方向导数 $d_1^3/(d_1^2+d_2^2)$ 存在但不是方向的线性函数，故原点不可微。

二次连续可微时，泰勒展开为 $f(x+h)=f(x)+\nabla f(x)^{\mathsf T}h+\tfrac12h^{\mathsf T}\nabla^2f(x)h+o(\|h\|_2^2)$。仿射复合 $g(z)=f(Az+b)$ 满足 $\nabla g=A^{\mathsf T}\nabla f$、$\nabla^2g=A^{\mathsf T}(\nabla^2f)A$。这些关系解释了降维和变量缩放对曲率的作用。

对称矩阵可谱分解为 $X=U\Lambda U^{\mathsf T}$，$U$ 正交；它半正定当且仅当全部特征值非负。正定时平方根和逆可通过对特征值分别开方、取倒数定义。谱范数为 $\|A\|_2=\max_{\|x\|_2=1}\|Ax\|_2$，等于最大奇异值。正定矩阵的 $\log\det X$ 沿方向 $H$ 的导数为 $\operatorname{tr}(X^{-1}H)$；逆矩阵的方向导数为 $-X^{-1}HX^{-1}$，后者可对恒等式 $XX^{-1}=I$ 求导得到。

\begin{example}[最小二乘梯度的维度与推导]
令 $A\in\mathbb R^{r\times n}$、$x\in\mathbb R^n$、$b\in\mathbb R^r$，求 $f(x)=\|Ax-b\|_2^2/2$ 的梯度和海塞矩阵。
\end{example}
\begin{solve}
设扰动为 $h\in\mathbb R^n$，展开得
\[
 f(x+h)-f(x)=(Ax-b)^{\mathsf T}Ah+\frac12h^{\mathsf T}A^{\mathsf T}Ah.
\]
把一次项写成 $[A^{\mathsf T}(Ax-b)]^{\mathsf T}h$，就读出梯度 $A^{\mathsf T}(Ax-b)\in\mathbb R^n$；二次项给出海塞矩阵 $A^{\mathsf T}A\in\mathbb R^{n\times n}$。残差有 $r$ 个分量，而梯度必须有 $n$ 个分量，维度检查可以立即排除把梯度误写为 $A(Ax-b)$ 的错误。
\end{solve}

实际计算矩阵微分时，先写出扰动的一次项，再按内积读出梯度，通常比逐元素求导更不易出错。正定矩阵的情形也一样：若 $X=U\Lambda U^{\mathsf T}\succ0$，则 $\log\det X=\sum_i\log\lambda_i$，它把各特征方向的尺度相加；当任一特征值趋于零时，$-\log\det X$ 趋于无穷，这解释了它作为半正定锥障碍函数的作用。

\section{概率工具与学习检查}

随机向量 $Z$ 的均值为 $\mu=\mathbb EZ$，协方差为 $\Sigma=\mathbb E[(Z-\mu)(Z-\mu)^{\mathsf T}]$。若二阶矩有限，对给定对称矩阵 $P$ 有
\[
 \mathbb E[Z^{\mathsf T}PZ]=\mu^{\mathsf T}P\mu+\operatorname{tr}(P\Sigma),
\]
可由 $Z=\mu+(Z-\mu)$ 展开并利用中心化项均值为零得到。独立性推出协方差为零，反向一般不成立。

交换参数微分与期望需要条件。例如在参数点的某邻域中，样本函数几乎处处可微，其各偏导绝对值受同一个可积随机变量控制，且函数在某参考参数处可积，则可用支配收敛定理交换微分与期望。计算似然或信息矩阵时，应验证这些条件或使用无需交换的有限和表达。

\begin{enumerate}
 \item 证明 $\|x\|_2\leq\|x\|_1\leq\sqrt n\|x\|_2$，并给出两端取等号的例子。
 \item 构造闭但不有界的可行集及连续目标，使有限下确界不取到。提示：$e^x$ 在实轴上。
 \item 对 $X\succ0$，由逆矩阵导数重新推导 $-\log\det X$ 的二阶方向导数。
\end{enumerate}

